\section{Introduction}

Quantum networks aim to distribute entanglement between distant nodes, unlocking
applications that range from provably secure key exchange to distributed sensing and
modular quantum computation~\cite{wehner2018quantum}. Because photon loss grows
exponentially with fibre length, entanglement cannot be sent directly over continental
distances; instead, short entangled links are joined at intermediate nodes through
\emph{entanglement swapping}, whereby a Bell-state measurement projects two independently
generated pairs onto a single longer-range entangled state~\cite{zukowski1993eventready}.
This stitching is not free: the secret-key rate achievable across a lossy link decays at
best linearly in the channel transmissivity, an information-theoretic ceiling that no
repeaterless protocol can surpass~\cite{wootters1998entanglement}.

Multi-hop routing over these quantum networks requires link-layer coordination and
path-selection algorithms that account for both decoherence and probabilistic
swap outcomes~\cite{pant2019routing,dahlberg2019linklayer}. Recent work has shown that
the order in which swaps are performed along a path significantly impacts the end-to-end
fidelity and success probability, making swap scheduling a first-class routing
consideration~\cite{bennett1993teleporting}.
